\subsection{Examples of valuations}

The examples of valuation rings given in § 1, no. 4 provide us with Examples 1 to 4 below:

Example (1) Every valuation on a finite field F is improper, since every element of \(F^{*}\) is a root unity.

Example (2) If K is a subfield of a field \(K'\) , the restriction to K of a valuation on \(K'\) is a valuation on K.

Example (3) Let \(k\) be a field and \(K = k((T))\) . The mapping \(v\) which maps every non-zero formal power series to its order (Algebra, Chapter IV, § 5, no. 7) is a valuation on \(K\) whose order group is \(Z\) and whose ring is \(k[[T]]\) . The associated place is the canonical homomorphism \(f: k[[T]] \to k\) extended to \(k((T))\) by setting \(f(u) = \infty\) if \(u \notin k[[T]]\) .

Example (4) Let A be a principal ideal domain, K its field of fractions and p an extremal element of A. For \(x \in K^{*}\) let \(v_{p}(x)\) denote the exponent of p in the decomposition of x into extremal elements (Algebra, Chapter VII, §1, no. 3, Theorem 2); it is immediately seen that \(v_{p}\) is a valuation whose order group is Z and whose ring is \(A_{Ap}\) . By Proposition 3 of §1, no. 4 we thus obtain, up to equivalence, all the valuations on K which are not improper and are positive on A. Taking A = Z we recover the p-adic valuations on Q (General Topology, Chapter IX, §3, no. 2); these valuations are, up to equivalence, the only valuations on Q which are not improper (§1, no. 4, Corollary 1 to Proposition 3). Taking A = k{[}X{]}, where k is a field, the non-improper valuations on k(X) whose restrictions to k are improper are (up to equivalence): on the one hand the valuations \(v_{P}\) where P runs through the set of irreducible monic polynomials of k{[}X{]} and on the other hand the valuation v defined by

\[
v (\mathrm{P} / \mathrm{Q}) = \deg (\mathrm{Q}) - \deg (\mathrm{P})
\]

for \(P \in k[X]\) and \(Q \in k[X]\) (§ 1, no. 4, Corollary 2 to Proposition 3); all these valuations obviously have Z as order group and their residue fields are monogenous algebraic extensions of k (Algebra, Chapter V, § 3, no. 1).

Example (5) The mapping \(P(X, Y) \mapsto P(T, e^{\mathrm{T}})\) of \(\mathbf{C}[X, Y]\) to \(\mathbf{C}((T))\) is injective (Functions of a real variable, Chapter IV, §2, Proposition 9) and therefore can be extended to an isomorphism of \(\mathbf{C}(X, Y)\) onto a subfield of \(\mathbf{C}((T))\) . The restriction to this subfield of the valuation on \(\mathbf{C}((T))\) defined in Example 3 defines a valuation on \(\mathbf{C}(X, Y)\) which is improper on C, whose order group is Z and whose residue field is C.

Proposition 4 of no. 2 allows us to construct a valuation whose order group and residue field are given:

Example (6) Let \(\Gamma\) be a totally ordered group and \(k\) a field. Let \(\Gamma_{+}\) be the monoid of positive elements of \(\Gamma\) and C the algebra of \(\Gamma_{+}\) over \(k\) . By definition, C has a basis \((x_{\alpha})_{\alpha \in \Gamma_{+}}\) over \(k\) whose multiplication table is \(x_{\alpha}x_{\beta} = x_{\alpha +\beta}\) . If \(x = \sum_{\alpha}a_{\alpha}x_{\alpha}\) is a non-zero element of C, we write \(v(x) = \inf_{a_{\alpha}\neq 0}(\alpha)\) and \(v(0) = +\infty\) ; it is immediately verified that the mapping \(v\) of C to \(\Gamma_{\infty}\) satisfies conditions \((\mathrm{VL}_{\mathrm{I}})\) and \((\mathrm{VL}_{\mathrm{II}})\) of no. 1 and that C is an integral domain. Let K be the field of fractions of C and \(w\) the valuation on K which extends \(v\) (Proposition 4, no. 2). As every element of \(\Gamma\) is the difference of two positive elements, \(w\) admits \(\Gamma\) as order group. Let A be the ring of \(w\) and m its maximal ideal; we shall show that A is the direct sum of m and k (identified with k.1), which will prove that the residue field of \(w\) is isomorphic to k. Clearly \(m \cap k = (0)\) . On the other hand, denoting by p the ideal of C generated by the \(x_{\alpha}\) where \(\alpha > 0\) , every element x of valuation 0 in K can be written in the form \((a + y)/(b + z)\) where \(a \in k^{*}\) , \(b \in k^{*}\) , \(y \in p\) and \(z \in p\) ; then

\[
x = a b ^ {- 1} + (b y - a z) b ^ {- 1} (b + z) ^ {- 1}
\]

whence \(w(x - ab^{-1}) > 0\) and \(x \equiv ab^{-1} \pmod{m}\) ; this shows our assertion.

If \(\Gamma = \mathbf{Z} \times \mathbf{Z}\) , then \(\mathbf{K} = k(\mathbf{X}, \mathbf{Y})\) and the above construction then provides valuations on \(k(\mathbf{X}, \mathbf{Y})\) which are improper on \(k\) , whose order group is \(\mathbf{Z} \times \mathbf{Z}\) and whose residue field is \(k\) . These valuations depend on the order structure chosen on \(\mathbf{Z} \times \mathbf{Z}\) . For example, \(\mathbf{Z} \times \mathbf{Z}\) can be given the lexicographic ordering. Or indeed, for an irrational number \(\alpha\) , \(\mathbf{Z} \times \mathbf{Z}\) may be identified with a subgroup of \(\mathbf{R}\) under the homomorphism \((m, n) \mapsto m + n\alpha\) (a homomorphism which is injective since \(\alpha\) is irrational) and given the ordering induced by that on \(\mathbf{R}\) .

Other constructions of valuations using Proposition 4 of no. 2 will be described in § 10.

